\documentclass[a4paper,12pt]{article}
\usepackage[T2A]{fontenc}
\usepackage[utf8]{inputenc}
\usepackage[russian]{babel}
\usepackage{amsmath,amssymb,mathtools}
\usepackage{icomma}
\usepackage[left=22mm,right=22mm,top=23mm,bottom=22mm]{geometry}
\usepackage{microtype}
\usepackage{tikz}
\usetikzlibrary{calc,positioning,arrows.meta,shapes.geometric}
\usepackage{tkz-euclide}
\usepackage{booktabs,tabularx}
\usepackage{enumitem}
\usepackage{fancyhdr}
\usepackage{setspace}
\usepackage{xcolor}
\usepackage[unicode,bookmarks=false,colorlinks=false,pdfborder={0 0 0}]{hyperref}
\definecolor{accent}{HTML}{008080}
\definecolor{darkaccent}{HTML}{004D4D}
\definecolor{textgray}{HTML}{333333}
\definecolor{quietgray}{HTML}{707878}
\definecolor{lightaccent}{HTML}{E8F4F1}
\color{textgray}
\onehalfspacing
\setlength{\parindent}{0pt}
\setlength{\parskip}{6pt}
\setcounter{secnumdepth}{1}
\renewcommand{\thesection}{\arabic{section}}
\usepackage{titlesec}
\titleformat{\section}{\large\color{darkaccent}}{\thesection.}{0.5em}{}
\titlespacing*{\section}{0pt}{17pt}{11pt}
\titleformat{\subsection}{\normalsize\color{accent}}{}{0pt}{}
\titlespacing*{\subsection}{0pt}{12pt}{5pt}
\pagestyle{fancy}
\fancyhf{}
\fancyhead[L]{\small Площади и диагонали трапеции}
\fancyhead[R]{\small София Кальней}
\fancyfoot[C]{\small\thepage}
\renewcommand{\headrulewidth}{0.35pt}
\renewcommand{\footrulewidth}{0pt}
\setlength{\headheight}{15pt}
\setlist[enumerate]{leftmargin=*,itemsep=7pt,topsep=4pt}
\newcommand{\idea}[1]{\textcolor{darkaccent}{\textbf{Ключевая идея.}}~#1}
\newcommand{\warn}[1]{\textcolor{accent}{\textbf{Обратите внимание.}}~#1}
\newcommand{\SA}{S_{AOD}}
\newcommand{\SB}{S_{BOC}}
\newcommand{\ST}{S_{ABCD}}

\begin{document}
\thispagestyle{empty}
\begin{center}
\vspace*{17mm}
{\Large\color{darkaccent} МАТЕМАТИКА}\par
\vspace{22mm}
{\huge Чек-лист}\par
\vspace{15mm}
{\LARGE\color{accent} Площади треугольников\par
в трапеции}\par
\vspace{17mm}
{\large Подобие, общая высота и диагонали}\par
\vspace{8mm}
{\normalsize Подготовка к ОГЭ по математике}\par
\vspace{12mm}
{\large София Кальней}\par
\vspace{9mm}
13 ноября 2026 г.\par
\vspace{11mm}
{\small Лёвин Артём Александрович\par
эксклюзивно для Софии Кальней}\par
\vspace{7mm}
\begin{tikzpicture}
\draw[accent, line width=0.9pt] (-3.2,0) .. controls (-2.2,0.42) and (-1.2,-0.43) .. (0,0)
.. controls (1.2,0.43) and (2.2,-0.42) .. (3.2,0);
\fill[darkaccent] (0,0) circle (1.35pt);
\end{tikzpicture}
\vspace*{8mm}
\end{center}
\clearpage

\section{Что изображено на чертеже}
Рассмотрим произвольную \textbf{выпуклую трапецию} $ABCD$, в которой основания $AD\parallel BC$. Её диагонали $AC$ и $BD$ пересекаются в точке $O$.

\begin{center}
\begin{tikzpicture}[scale=0.86]
\tkzDefPoint(0,0){A}
\tkzDefPoint(8,0){D}
\tkzDefPoint(1.2,3.2){B}
\tkzDefPoint(5.9,3.2){C}
\tkzInterLL(A,C)(B,D)\tkzGetPoint{O}
\fill[accent!13] (A)--(O)--(D)--cycle;
\fill[accent!6] (B)--(O)--(C)--cycle;
\fill[gray!7] (A)--(B)--(O)--cycle;
\fill[gray!7] (D)--(C)--(O)--cycle;
\tkzDrawPolygon[line width=1.0pt](A,B,C,D)
\tkzDrawSegments[line width=1pt](A,C B,D)
\tkzDrawPoints[size=3](A,B,C,D,O)
\tkzLabelPoint[below left](A){$A$}
\tkzLabelPoint[above left](B){$B$}
\tkzLabelPoint[above right](C){$C$}
\tkzLabelPoint[below right](D){$D$}
\tkzLabelPoint[above](O){$O$}
\node[darkaccent] at (4.05,0.67) {$S_1$};
\node[darkaccent] at (3.5,2.86) {$S_2$};
\node[quietgray] at (1.8,1.66) {$T_1$};
\node[quietgray] at (5.4,1.64) {$T_2$};
\end{tikzpicture}
\end{center}

Точка $O$ разбивает трапецию на четыре треугольника. Обозначим
\[
\SA=S_1,\qquad \SB=S_2,\qquad S_{AOB}=T_1,\qquad S_{COD}=T_2.
\]
Их площади складываются без перекрытия:
\[
\boxed{\ST=S_1+S_2+T_1+T_2.}
\]

\idea{Чтобы найти площадь всей трапеции, достаточно определить площади двух боковых треугольников.}

\section{Какие треугольники подобны}
Из параллельности оснований и вертикальных углов следует
\[
\triangle AOD\sim\triangle COB.
\]
Соответствующие стороны: $AO\leftrightarrow CO$, $DO\leftrightarrow BO$, $AD\leftrightarrow CB$. Следовательно,
\[
\frac{AO}{CO}=\frac{DO}{BO}=\frac{AD}{BC}=k>0,
\qquad \frac{\SA}{\SB}=k^2.
\]
Отсюда, при $S_1,S_2>0$,
\[
\boxed{k=\sqrt{\frac{S_1}{S_2}}.}
\]
\warn{Отношение площадей подобных треугольников равно \emph{квадрату} коэффициента подобия.}
\clearpage

\section{Общая высота и отношение площадей}
Треугольники $AOB$ и $BOC$ имеют основания $AO$ и $OC$ на одной прямой $AC$. Высота из точки $B$ для них \textbf{общая}.

\begin{center}
\begin{tikzpicture}[scale=0.72]
\tkzDefPoint(0,0){A}
\tkzDefPoint(8,0){D}
\tkzDefPoint(1.2,3.2){B}
\tkzDefPoint(5.9,3.2){C}
\tkzInterLL(A,C)(B,D)\tkzGetPoint{O}
\coordinate (H) at ($(A)!(B)!(C)$);
\fill[gray!9] (A)--(B)--(O)--cycle;
\fill[accent!10] (B)--(O)--(C)--cycle;
\tkzDrawPolygon[line width=0.8pt](A,B,C,D)
\tkzDrawSegments[line width=0.9pt](A,C B,D)
\draw[accent,densely dashed,line width=0.9pt] (B)--(H);
\tkzDrawPoints[size=3](A,B,C,D,O)
\tkzLabelPoint[below left](A){$A$}
\tkzLabelPoint[above left](B){$B$}
\tkzLabelPoint[above right](C){$C$}
\tkzLabelPoint[below right](D){$D$}
\tkzLabelPoint[above right](O){$O$}
\node[accent,left] at (H){$h$};
\end{tikzpicture}
\end{center}

По формуле площади треугольника получаем
\[
\begin{aligned}
S_{AOB}&=\frac12 AO\cdot h,\\
S_{BOC}&=\frac12 OC\cdot h.
\end{aligned}
\qquad\Longrightarrow\qquad
\frac{S_{AOB}}{S_{BOC}}=\frac{AO}{OC}=k.
\]
Подставляем $S_{BOC}=S_2$ и найденное по подобию отношение сторон:
\[
S_{AOB}=S_2\sqrt{\frac{S_1}{S_2}}
=\boxed{\sqrt{S_1S_2}}.
\]

Аналогично, у треугольников $COD$ и $BOC$ общая высота из точки $C$ к прямой $BD$:
\[
\frac{S_{COD}}{S_{BOC}}=\frac{DO}{BO}=k,
\qquad S_{COD}=\boxed{\sqrt{S_1S_2}}.
\]

\idea{Два боковых треугольника равновелики: $S_{AOB}=S_{COD}$. Их формы и длины сторон могут различаться.}

\subsection{Ещё одно доказательство равенства площадей}
Треугольники $ABD$ и $ACD$ имеют общее основание $AD$ и одинаковую высоту (вершины $B$ и $C$ лежат на прямой $BC\parallel AD$). Значит,
\[
S_{ABD}=S_{ACD}\;\Longrightarrow\;
S_1+S_{AOB}=S_1+S_{COD}\;\Longrightarrow\;S_{AOB}=S_{COD}.
\]
\clearpage

\section{Главная формула площади трапеции}
Сложим площади всех четырёх треугольников:
\[
\ST=S_1+\sqrt{S_1S_2}+\sqrt{S_1S_2}+S_2.
\]
Поскольку $S_1,S_2>0$, имеем $\sqrt{S_1S_2}=\sqrt{S_1}\sqrt{S_2}$. Поэтому
\[
\begin{aligned}
\ST&=(\sqrt{S_1})^2+2\sqrt{S_1}\sqrt{S_2}+(\sqrt{S_2})^2\\
&=\boxed{\bigl(\sqrt{S_1}+\sqrt{S_2}\bigr)^2}.
\end{aligned}
\]
Это \textbf{квадрат суммы}: $a^2+2ab+b^2=(a+b)^2$.

\subsection{Пример с занятия: $S_1=25$, $S_2=9$}
\textbf{Шаг 1.} Определим коэффициент подобия:
\[
k=\sqrt{\frac{25}{9}}=\frac{5}{3}.
\]
\textbf{Шаг 2.} Найдём площадь бокового треугольника, используя общую высоту:
\[
S_{AOB}=\frac{AO}{CO}\cdot S_{BOC}=\frac53\cdot 9=15.
\]
\textbf{Шаг 3.} Для второго бокового треугольника получаем
\[
S_{COD}=\frac{DO}{BO}\cdot S_{BOC}=\frac53\cdot 9=15.
\]
\textbf{Шаг 4.} Складываем:
\[
\boxed{\ST=25+15+15+9=64.}
\]
\textbf{Проверка по короткой формуле:}
\[
(\sqrt{25}+\sqrt9)^2=(5+3)^2=64.
\]

\subsection{Обратный ход: если известна площадь бокового треугольника}
Из $T=\sqrt{S_1S_2}$ следует $T^2=S_1S_2$. Поэтому при известных $T$ и $S_1>0$
\[
\boxed{S_2=\frac{T^2}{S_1}}.
\]
Например, если $T=15$ и $S_1=25$, то $S_2=225/25=9$.
\clearpage

\section{Что важно не перепутать}
\begin{enumerate}[label=\arabic*.]
\item \textbf{Какие треугольники подобны?} \quad $\triangle AOD$ и $\triangle COB$ (треугольники при основаниях). Для пары $AOB$ и $COD$ доказано равенство \emph{площадей}; подобие этой пары для произвольной трапеции не требуется.
\item \textbf{Сторона и площадь.} \quad Если отношение соответствующих сторон равно $5:3$, то отношение площадей подобных треугольников равно $25:9$.
\item \textbf{Порядок в дроби.} \quad При $AO/CO$ наверху ставьте площадь треугольника с отрезком $AO$, то есть $S_{AOD}$; внизу --- с $CO$, то есть $S_{BOC}$.
\item \textbf{Общая высота.} \quad Отношение площадей равно отношению оснований только для треугольников, у которых высота к этим основаниям одинакова.
\item \textbf{Корень и квадрат.} \quad $\sqrt{25/9}=5/3$, а $\sqrt{25\cdot9}=15$. Перенос положительного множителя под корень требует возведения его в квадрат.
\item \textbf{Обязательное доказательство.} \quad В развёрнутом решении покажите подобие, связь площадей с общей высотой и вычисление двух боковых площадей, после чего складывайте четыре части.
\end{enumerate}

\subsection{Краткая памятка}
\begin{center}
\renewcommand{\arraystretch}{1.45}
\begin{tabularx}{\linewidth}{@{}p{.40\linewidth}X@{}}
\toprule
Что известно или замечено & Что делать \\
\midrule
Параллельные основания & Найти подобные треугольники $AOD$ и $COB$ \\
Площади $S_1$, $S_2$ при основаниях & Вычислить $k=\sqrt{S_1/S_2}$ \\
Общая высота двух треугольников & Сравнить их площади через основания \\
Два боковых треугольника & Их площади одинаковы: $\sqrt{S_1S_2}$ \\
Нужна вся трапеция & Использовать $(\sqrt{S_1}+\sqrt{S_2})^2$ \\
\bottomrule
\end{tabularx}
\end{center}
\clearpage

\thispagestyle{fancy}
\begin{center}
{\Large\color{darkaccent} Алгоритм решения задачи о площадях в трапеции}\par
\vspace{10mm}
\end{center}
\begin{center}
\begin{tikzpicture}[
  node distance=8mm and 13mm,
  start/.style={ellipse,draw=darkaccent,line width=.8pt,minimum height=10mm,text width=32mm,align=center},
  action/.style={rectangle,rounded corners=2mm,draw=accent,line width=.7pt,text width=88mm,minimum height=11mm,align=center,inner sep=5pt},
  test/.style={diamond,aspect=3.3,draw=darkaccent,line width=.8pt,text width=60mm,align=center,inner xsep=0pt,inner ysep=4pt},
  arrow/.style={-{Stealth[length=2.3mm]},line width=.8pt,darkaccent}
]
\node[start] (begin) {Начало};
\node[action,below=of begin] (select) {Обозначить $S_1=S_{AOD}$, $S_2=S_{BOC}$};
\node[test,below=11mm of select] (known) {Известны обе площади $S_1$ и $S_2$?};
\node[action,below=13mm of known] (similar) {По подобию найти $k=\sqrt{S_1/S_2}$};
\node[action,below=of similar] (sides) {По общей высоте найти боковые площади: $T_1=T_2=\sqrt{S_1S_2}$};
\node[action,below=of sides] (sum) {Сложить четыре площади: $S=S_1+T_1+T_2+S_2$};
\node[action,below=of sum] (finalf) {$\displaystyle S=(\sqrt{S_1}+\sqrt{S_2})^2$};
\node[start,below=of finalf] (end) {Конец};
\node[action,right=19mm of known,text width=36mm] (recover) {Найти недостающую площадь по данным условия};
\draw[arrow] (begin)--(select);
\draw[arrow] (select)--(known);
\draw[arrow] (known)--node[left,font=\small]{да}(similar);
\draw[arrow] (known.east)--node[above,font=\small]{нет}(recover.west);
\draw[arrow] (recover.north) |- ([yshift=5mm]known.north) -- (known.north);
\draw[arrow] (similar)--(sides);
\draw[arrow] (sides)--(sum);
\draw[arrow] (sum)--(finalf);
\draw[arrow] (finalf)--(end);
\end{tikzpicture}
\end{center}
\clearpage

\section*{Тренировочный блок}
\textcolor{accent}{\textbf{Путь А --- 10 заданий}}\par
Во всех заданиях $ABCD$ --- трапеция, $AD\parallel BC$, $O$ --- точка пересечения диагоналей. Если требуется развёрнутое решение, обоснуйте все использованные соотношения.

\subsection{Базовый уровень}
\begin{enumerate}[label=\textbf{\arabic*.},itemsep=9pt]
\item Основания трапеции относятся как $AD:BC=2:3$. Найдите отношение $S_{AOD}:S_{BOC}$.
\item Даны $S_{AOD}=25$ и $S_{BOC}=9$. Найдите $AO:CO$.
\item Даны $S_{AOD}=9$ и $S_{BOC}=16$. Найдите площади обоих боковых треугольников $AOB$ и $COD$.
\item Даны $S_{AOD}=4$ и $S_{BOC}=9$. Найдите площадь трапеции.
\end{enumerate}

\subsection{Средний уровень}
\begin{enumerate}[label=\textbf{\arabic*.},start=5,itemsep=9pt]
\item Даны $S_{AOD}=25$ и $S_{BOC}=9$. Подробно докажите равенство площадей $AOB$ и $COD$, затем найдите площадь трапеции.
\item Даны $S_{AOD}=12$ и $S_{BOC}=27$. Найдите площадь трапеции.
\item Даны $S_{AOD}=25$ и $S_{AOB}=15$. Найдите $S_{BOC}$ и площадь всей трапеции.
\end{enumerate}

\subsection{Повышенный уровень}
\begin{enumerate}[label=\textbf{\arabic*.},start=8,itemsep=9pt]
\item Даны $S_{BOC}=16$ и $S_{COD}=24$. Найдите $S_{AOD}$ и площадь трапеции.
\item Основания относятся как $AD:BC=3:2$, а $S_{AOD}=27$. Найдите площадь трапеции.
\item Площадь трапеции равна $121$, а $S_{AOD}=49$. Найдите $S_{BOC}$ и площади двух боковых треугольников.
\end{enumerate}
\clearpage

\section*{Ответы к тренировочному блоку}
\vspace{6mm}
\renewcommand{\arraystretch}{1.60}
\begin{center}
\begin{tabularx}{\linewidth}{@{}p{14mm}X@{}}
\toprule
\textbf{№} & \textbf{Ответ} \\
\midrule
1 & $4:9$ \\
2 & $5:3$ \\
3 & $S_{AOB}=12$, $S_{COD}=12$ \\
4 & $25$ \\
5 & $S_{AOB}=S_{COD}=15$, площадь трапеции $64$ \\
6 & $75$ \\
7 & $S_{BOC}=9$, площадь трапеции $64$ \\
8 & $S_{AOD}=36$, площадь трапеции $100$ \\
9 & $75$ \\
10 & $S_{BOC}=16$, $S_{AOB}=S_{COD}=28$ \\
\bottomrule
\end{tabularx}
\end{center}
\vspace{12mm}
\subsection*{Итоговая памятка}
\begin{align*}
\triangle AOD&\sim\triangle COB,\\[2mm]
\frac{AO}{CO}&=\frac{DO}{BO}=\frac{AD}{BC}=\sqrt{\frac{S_1}{S_2}},\\[2mm]
S_{AOB}&=S_{COD}=\sqrt{S_1S_2},\\[2mm]
\boxed{\ST=(\sqrt{S_1}+\sqrt{S_2})^2}
\end{align*}
\vfill
\begin{center}
\small\color{quietgray} Лёвин Артём Александрович --- эксклюзивно для Софии Кальней
\end{center}
\end{document}
